\section{用例}
\subsection{フォント}

\subsubsection{文字修飾}
\begin{itemize}
\item 普通の文字(normal)
\item \emph{強調(emph)}
\item \sout{打ち消し線}
\item \textsuperscript{上付き}
\end{itemize}

\subsubsection{書体}

\begin{itemize}
\item \textrm{ローマン体}
\item \textsf{サンセリフ体}
\item \texttt{タイプライター体}
\item \textmd{デフォルト体}
\item \textbf{ボールド}
\item \textup{直立体}
\item \textit{イタリック}
\item \textsl{スラント}
\item \textsc{スモールキャピタル}
\end{itemize}

\subsubsection{フォントサイズ}
\def\sampletext{Windowsでコンピュータの世界が広がります。}
\begin{itemize}
  \normalsize\item\tiny{\sampletext(tiny)}
  \normalsize\item\scriptsize{\sampletext(scriptsize)}
  \normalsize\item\footnotesize{\sampletext(footnotesize)}
  \normalsize\item\small{\sampletext(small)}
  \normalsize\item\normalsize{\sampletext(normalsize)}
  \normalsize\item\large{\sampletext(large)}
  \normalsize\item\Large{\sampletext(Large)}
  \normalsize\item\LARGE{\sampletext(LARGE)}
  \normalsize\item\huge{\sampletext(huge)}
  \normalsize\item\Huge{\sampletext(Huge)}
\end{itemize}
\normalsize

\subsection{リンク}
\url{http://localhost/}

\href{http://www.example.com/}{example.com}

\subsection{quote環境}

\begin{quote}
吾輩は猫である。名前はまだ無い。

どこで生れたかとんと見当がつかぬ。
\end{quote}

\subsection{quotation環境}

\begin{quotation}
吾輩は猫である。名前はまだ無い。

どこで生れたかとんと見当がつかぬ。
\end{quotation}

\subsection{verbatim環境}

半角80文字がちょうど収まるサイズにしています。
コマンドラインやソースコードを示す場合はverbatim環境よりもlistingsを推奨します

\begin{verbatim}
01234567890123456789012345678901234567890123456789012345678901234567890123456789
０１２３４５６７８９０１２３４５６７８９０１２３４５６７８９０１２３４５６７８９
\end{verbatim}

フォントが小さくなって嫌だ、という方のために60文字がちょうど収まるverbatim60とverbatim40を用意しています。

\subsection{verbatim60}

\begin{verbatim60}
012345678901234567890123456789012345678901234567890123456789
０１２３４５６７８９０１２３４５６７８９０１２３４５６７８９
\end{verbatim60}

\subsection{verbatim40}

\begin{verbatim40}
0123456789012345678901234567890123456789
０１２３４５６７８９０１２３４５６７８９
\end{verbatim40}

\subsection{コードリスティング}

\subsubsection{コマンドライン}
\begin{lstlisting}[language=command]
# uname --machine
x86_64
\end{lstlisting}

\subsubsection{C}
\begin{lstlisting}[language=C]
#include <stdio.h>

int main(int argc, char *argv[]){
    printf("Hello World!");
    return 0;
}
\end{lstlisting}

\subsubsection{Python}
\begin{lstlisting}[language=python]
def hello():
    print('Hello World!')

if __name__ == '__main__':
    hello()
\end{lstlisting}

\subsubsection{80文字テスト}
\begin{lstlisting}
01234567890123456789012345678901234567890123456789012345678901234567890123456789
\end{lstlisting}

\subsection{リスト}

\subsubsection{列挙リスト}
\begin{itemize}
\item りんご
\item みかん
\item バナナ
\end{itemize}

\subsubsection{順序リスト}
\begin{enumerate}
\item 睦月
\item 如月
\item 弥生
\end{enumerate}

\subsubsection{定義リスト}
\begin{description}
\item[Monday] 月曜日
\item[Tuesday] 火曜日
\item[Wednesday] 水曜日
\end{description}

\subsection{テーブル}

\subsubsection{普通のテーブル}
\begin{table}[htbp]
\begin{center}
\begin{tabular}{|l|c|r|}
\hline
セル1 & セル2 & セル3 \\ \hline
セル4 & セル5 & セル6 \\ \hline
セル7 & セル8 & セル9 \\ \hline
\end{tabular}
\caption{キャプション}
\end{center}
\end{table}

\subsubsection{longtable}
\begin{longtable}[c]{@{}crr@{}}
\toprule\addlinespace
果物 & 値段 & 個数
\\\addlinespace
\midrule\endhead
バナナ & 134円 & 10本
\\\addlinespace
オレンジ & 210円 & 10個
\\\addlinespace
\bottomrule
\end{longtable}

\subsection{水平線}
\begin{center}\rule{3in}{0.4pt}\end{center}

\subsection{画像}
\includegraphics[width=\textwidth]{latex.eps}

\clearpage
\subsection{数式}
\begin{enumerate}

\item 質量とエネルギーの等価性
\begin{equation}
E = mc ^2
\end{equation}

\item 波動方程式
\begin{equation}
\frac{\partial^2 z}{\partial t^2}=
c^2 (\frac{\partial^2 z}{\partial x^2}+\frac{\partial^2 z}{\partial y^2})-
\mu \frac{\partial z}{\partial t}
\end{equation}

\begin{equation}
\sum_{x=1}^4 x^2 = 30
\end{equation}

\item オイラーの公式

文章中の $e^{i\theta}=\cos\theta+i\sin\theta$ オイラーの公式。

\end{enumerate}
